% Generated draft: requires mathematical and source review.
\documentclass[11pt]{article}
\usepackage[margin=25mm]{geometry}
\usepackage{fontspec}
\setmainfont{Latin Modern Roman}
\usepackage{amsmath,amssymb,mathtools,unicode-math}
\setmathfont{Latin Modern Math}
\usepackage{enumitem,needspace,xurl,hyperref}
\setlength{\parindent}{0pt}
\setlength{\parskip}{5pt}
\setlength{\emergencystretch}{3em}
\newcommand{\N}{\mathbb{N}}
\newcommand{\Z}{\mathbb{Z}}
\newcommand{\Q}{\mathbb{Q}}
\newcommand{\R}{\mathbb{R}}
\newcommand{\C}{\mathbb{C}}
\newcommand{\F}{\mathbb{F}}
\newcommand{\A}{\mathbb{A}}
\newcommand{\PP}{\mathbb{P}}
\newcommand{\E}{\mathbb{E}}
\newcommand{\eps}{\varepsilon}
\newcommand{\dd}{\,\mathrm{d}}
\begin{document}
\section*{UnsolvedMath ID: 30002532}
\textbf{Source problem code:} catalogue entry
\subsection{Definitions and mathematical statement}
Let \(X\) be a smooth projective variety of dimension \(n\) over a field \(k\). This notation is conditional: the supplied record does not specify the field, its characteristic, or any further hypotheses on \(X\). A codimension-\(p\) algebraic cycle on \(X\) is a finite formal sum \(\alpha=\sum_i m_i[V_i]\), where \(m_i\in\Z\) and each \(V_i\subseteq X\) is an integral closed subvariety of codimension \(p\). Such a cycle is effective if every \(m_i\geq 0\). The record does not specify whether the intended objects are integral subvarieties, effective cycles, arbitrary cycles, or classes modulo an equivalence relation.

The supplied research question is: \emph{Which natural higher-dimensional generalizations of the Bounded Negativity Conjecture hold for cycles on smooth projective varieties?} It asks for the formulation and investigation of suitable generalizations; it does not supply a particular quantified conjecture.

The intended higher-dimensional negativity invariant and the proposed bound are unresolved. For example, an intersection product of two codimension-\(p\) cycle classes has codimension \(2p\), so its degree gives a self-intersection number directly only when \(2p=n\). Other dimensions require a different construction or additional intersection data, neither of which is specified. The record also leaves unspecified the permitted cycle dimensions, any positivity or normalization conditions, and the data on which a bound may depend. Consequently, this draft preserves the research question without selecting a particular generalization or asserting an inequality.
\subsection{Short English statement}
Determine suitable formulations of bounded negativity for algebraic cycles on smooth projective varieties of higher dimension, and determine which formulations hold. The supplied question leaves the field, cycle class, numerical invariant, and quantifiers for the proposed bounds unspecified, so it does not determine a single precise conjecture.
\subsection{Sources}
\begin{itemize}
\item Local UnsolvedMath registry, permanent ID \(30002532\), title \emph{Higher-Dimensional Bounded Negativity for Algebraic Cycles}, source problem code \texttt{catalogue entry}, in the Oberwolfach Reports Open Problems collection. The supplied proposed year is 2014; no individual proposer is identified.
\item Supplied source citation: \emph{Mini-Workshop: Negative Curves on Algebraic Surfaces} (2014), Oberwolfach Reports, DOI 10.4171/owr/2014/10, pp.\ 555--593. \url{https://doi.org/10.4171/owr/2014/10}. The supplied background describes these reports as documenting mathematical workshops and their open-problem discussions; no exact problem page is supplied.
\item Additional supplied literature source: Bauer et al., 2023. \url{https://doi.org/10.1515/advgeom-2023-0027}. No title, complete author list, or specific supporting passage is supplied.
\item Literature assessment dated exactly \texttt{2026-08-21}: the supplied curation labels the problem open and states that higher-dimensional analogues of bounded negativity have no general established formulation or theorem resolving this research question. It lists the Oberwolfach report and Bauer et al.\ as sources checked. This is supplied dated curation derived from the OWR curation, not an independently verified assessment; the record explicitly requests independent verification of time-sensitive status and cited claims.
\item The registry labels the statement \texttt{corrected\_verified} and reports a manual semantic, typographic, and LaTeX review of its cleanup. These are supplied editorial qualifications, not independent verification of a particular mathematical formulation or its current status.
\end{itemize}
\end{document}
